%---
% title: "Chemokinetic Optimization: The Information Barrier"
% date: 2026-05-10
% work_order: 9
% permalink: /posts/2026/05/chemokinesis-curse-of-dimensionality/
% tags:
%   - optimization
%   - probability
%   - high-dimensional-statistics
%---
\documentclass[11pt]{article}
\usepackage{publication-web}
\begin{document}

Consider optimization when the dynamics supply an isotropic random unit direction \(u_t\) and the controller chooses only a nonnegative speed:

\[
\theta_{t+1}=\theta_t+s_tu_t,\qquad s_t\geq0.
\]

The ray cannot be reversed or probed before the move; feedback is a scalar loss. Locally,

\[
L(\theta+su)-L(\theta)=qs+\tfrac12hs^2+o(s^2),
\qquad q=g^\top u,\quad h=u^\top Hu.
\]

The problem separates into exploiting a ray when \(q\) is known and predicting \(q\) causally when it is not.

\section{Oracle Ray-CK and effective dimension}\label{oracle-ray-ck}

If \(q\) and \(h>0\) are known before moving, the optimal admissible action is

\[
s^*(u)=\frac{(-q)_+}{h}.
\]

It stops on uphill rays and never reverses them. For a quadratic \(L(x)=\tfrac12x^\top Hx\), the best scalar-multiplier version of this update loses a factor \(2d_{\mathrm{eff}}\) against exact gradient line search, where

\[
d_{\mathrm{eff}}(g,H)
=\frac{d}{d+2}\left(\frac{\operatorname{tr}H}{\lambda_g}+2\right),
\qquad
\lambda_g=\frac{g^\top Hg}{\lVert g\rVert^2},
\qquad
\frac{\Delta^*_{\mathrm{Ray}}}{\Delta^*_{\mathrm{GD}}}
=\frac{1}{2d_{\mathrm{eff}}}.
\]

If \(H=\lambda P\) for a rank-\(r\) projector and \(g\in\operatorname{range}(P)\), then \(d_{\mathrm{eff}}=d(r+2)/(d+2)\to r+2\). Flat nuisance coordinates therefore need not increase oracle cost.

\section{The causal information barrier}\label{causal-information-barrier}

Without a same-ray probe, the relevant quantities are the posterior directional slope and curvature,

\[
\mu_t(u_t)=\mathbb E[u_t^\top g_t\mid\mathcal F_t^-,u_t],
\qquad
\kappa_t(u_t)=\mathbb E[u_t^\top H_tu_t\mid\mathcal F_t^-,u_t].
\]

The Bayes action under the local quadratic model is \(s_t^*=(-\mu_t)_+/\kappa_t\). Directly predicting the ray response cannot beat the information available for estimating \(\mu_t\).

For \(g\sim\mathcal N(0,\sigma^2I_d)\), \(M\leq d\) noiseless random measurements \(y_i=u_i^\top g\) reveal only the projection of \(g\) onto their span \(S\). With constant ray curvature, the posterior predictor retains exactly

\[
\frac{\mathbb E[(u^\top P_Sg)^2]}{\mathbb E[(u^\top g)^2]}
=\frac{M}{d}
\]

of oracle-Ray progress. Parallel measurements remove latency, not this orientation cost. Standard zeroth-order methods avoid the barrier by probing and retaining the current direction; SPSA and Random Pursuit are therefore solving an easier information problem (\href{https://doi.org/10.1109/9.119632}{Spall 1992}; \href{https://doi.org/10.1137/110853613}{Stich et al.~2013}).

\section{Projected oracle and operational optimizer}\label{projected-and-operational}

After a warmup reveals \(P_Sg\), the projected oracle uses

\[
s_S^*(u)=\frac{(-u^\top P_Sg)_+}{u^\top Hu}.
\]

Under locally isotropic ray curvature, its expected progress fraction is \(\lVert P_Sg\rVert^2/\lVert g\rVert^2\), averaging to \(m/d\) for a random \(m\)-dimensional \(S\). Partial information remains useful because a fresh ray need not lie inside \(S\).

A lean operational version obtains parallel finite-difference slopes \(y_i=[L(\theta_0+\delta u_i)-L(\theta_0)]/\delta\), stacks the directions as rows of \(V\), and fits

\[
\widehat g=V^\top(VV^\top+\lambda I)^{-1}y.
\]

It then applies \(\theta_{t+1}=\theta_t+a_t(-u_t^\top\widehat g_t)_+u_t\), periodically adding displacement--loss chords and refitting. One global ray scale replaces an unstable online Hessian model.

\section{MNIST check}\label{mnist-check}

On MNIST, the operational method reached \(2.05\) negative log-likelihood versus \(1.97\) for gradient descent. Matching a GD trajectory required about \(30\) loss evaluations per GD step for oracle Ray-CK and about \(200\) for the operational optimizer. The small endpoint gap but large evaluation gap isolates the problem: ray exploitation is governed by \(d_{\mathrm{eff}}\), while causal orientation acquisition dominates the implementable method.

\section{Derivations}\label{derivations}

For \(x^+=x+\alpha(-g^\top U)_+U\) with \(U\) uniform on the sphere,

\[
\mathbb E[L(x^+)-L(x)\mid x]
=-\frac{\alpha}{2d}\lVert g\rVert^2
+\frac{\alpha^2}{4d(d+2)}
\left(\operatorname{tr}(H)\lVert g\rVert^2+2g^\top Hg\right).
\]

Minimizing in \(\alpha\) and dividing by the exact gradient-line-search decrease \(\lVert g\rVert^4/(2g^\top Hg)\) gives the \(2d_{\mathrm{eff}}\) ratio.

Conditioning an isotropic Gaussian on \(M\) exact projections leaves an independent isotropic component in \(S^\perp\), so \(\widehat g=P_Sg\). Since \(\mathbb E_u[(u^\top a)^2]=\lVert a\rVert^2/d\) and \(\mathbb E\lVert P_Sg\rVert^2=(M/d)\mathbb E\lVert g\rVert^2\), the causal and projected-oracle fractions follow.

\end{document}
